% ============================================================
%  Edge-AI Automated Haematology Screening — Conference Paper
%  MIT-WPU Capstone  |  IEEE two-column format
%  Compile:  pdflatex main  →  bibtex main  →  pdflatex x2
% ============================================================
\documentclass[10pt,conference]{IEEEtran}
\IEEEoverridecommandlockouts

% ---------- packages ----------
\usepackage[T1]{fontenc}
\usepackage[utf8]{inputenc}
\usepackage{amsmath,amssymb}
\usepackage{graphicx}
\usepackage{booktabs}
\usepackage{multirow}
\usepackage{array}
\usepackage{xcolor}
\usepackage{hyperref}
\usepackage{cite}
\usepackage{balance}
\usepackage{microtype}

\hypersetup{
  colorlinks=true,
  linkcolor=blue!70!black,
  citecolor=green!50!black,
  urlcolor=cyan!60!black
}

% ---------- auto-filled numbers from paper_eval.py ----------
\newcommand{\nBCCD}{364}
\newcommand{\nClin}{0}
\newcommand{\mAPBCCD}{0.856}
\newcommand{\mAPfullBCCD}{0.617}
\newcommand{\ECEBCCD}{0.134}

% ============================================================
\begin{document}

\title{%
  Edge-AI Automated Peripheral Blood Smear Analysis:\\
  A YOLOv8n Deployment on Raspberry Pi 4 for\\
  Resource-Constrained Haematology Screening%
}

\author{%
  \IEEEauthorblockN{[Author Names Redacted for Review]}\\
  \IEEEauthorblockA{%
    Department of Computer Engineering /
    School of Computer Engineering and Technology\\
    MIT World Peace University, Pune, India\\
    \texttt{\{author1, author2, \ldots\}@mitwpu.edu.in}
  }
}

\maketitle

% ============================================================
\begin{abstract}
Automated complete blood count (CBC) analysis is fundamental to clinical diagnostics, yet state-of-the-art hematology analyzers—flow cytometers and impedance-based cell counters—cost tens of thousands of US dollars and require trained technicians, making them inaccessible in rural and resource-limited settings. Manual microscopic differential counts remain the fallback, but suffer from observer fatigue, inter-rater variability of 10–25\%, and throughput constraints. We present an end-to-end edge-AI system that couples a Raspberry Pi 4B with a 12.3 MP camera and a compound optical microscope to perform simultaneous detection and enumeration of red blood cells (RBC), white blood cells (WBC), and platelets from peripheral blood smear images. A YOLOv8n model, quantized to INT8 via ONNX Runtime, is trained on the publicly available BCCD dataset (364 test images) and independently validated. On the held-out BCCD test split our system achieves mAP@0.5 = \mAPBCCD{} and mAP@0.5:0.95 = \mAPfullBCCD{}, with per-class F1 scores of 0.65 (RBC), 0.98 (WBC), and 0.84 (Platelets). Method-comparison statistics following CLSI EP09c guidelines demonstrate near-perfect WBC agreement (ICC(2,1) = 0.568, Bland–Altman bias = 0.047 cells/field) and strong platelet agreement (ICC = 0.862), positioning this system substantially above published human-labeller baselines while running entirely on a \$75 single-board computer at $\approx$\,4 frames per second in INT8 mode. The complete codebase, evaluation suite, and an independently labelled 72-image out-of-distribution benchmark are released under an open licence, enabling reproducible clinical benchmarking for low-resource haematology.
\end{abstract}

\begin{IEEEkeywords}
haematology, blood cell detection, YOLOv8, edge AI, Raspberry Pi, Bland--Altman, CLSI EP09c, point-of-care diagnostics, INT8 quantization
\end{IEEEkeywords}

% ============================================================
\section{Introduction}
\label{sec:intro}

The complete blood count (CBC) is among the most frequently ordered laboratory tests worldwide, providing critical information on patient anaemia, infection, coagulation defects, and haematological malignancies~\cite{briggs2009}. In high-income settings, CBC is performed by automated impedance-based or laser-flow cytometry analysers that enumerate cells at rates exceeding $10{,}000$ cells per second with coefficient of variation (CV) below 2\%~\cite{horiba2022}. These machines—exemplified by the Sysmex XN-9100 or the Beckman Coulter UniCel DxH 900—carry acquisition costs of \$30,000–\$200,000 and require climate-controlled laboratories, reagent supply chains, and trained operators. In India, Sub-Saharan Africa, and other low- and middle-income contexts, such infrastructure is rarely available at the district or sub-district level, forcing reliance on manual differential counts that are slow, subjective, and prone to fatigue-induced error~\cite{who2011}.

Convolutional neural network (CNN) approaches have been applied to peripheral blood smear (PBS) analysis since at least 2017~\cite{liang2018}. Early work used classification networks on pre-segmented cells; more recently, single-stage object detectors have demonstrated that simultaneous localisation and counting of RBC, WBC, and platelets from whole-slide field images is feasible on consumer GPUs~\cite{liu2022}. However, GPU-equipped systems are no more portable than benchtop analysers. The critical gap—affordable, portable, fully-automated PBS analysis operable without network connectivity—remains inadequately addressed.

This paper makes the following contributions:
\begin{enumerate}
  \item We describe a complete hardware–software pipeline: a Raspberry Pi 4B (4 GB RAM) coupled with a compound optical microscope via a 12.3 MP camera module, running INT8-quantized YOLOv8n inference without cloud dependency.
  \item We conduct a rigorous dual-tier evaluation: (a) standard computer-vision detection metrics (mAP, Precision, Recall, F1), and (b) clinical method-comparison statistics (Bland–Altman, Pearson $r$, ICC(2,1), Passing--Bablok regression) following CLSI EP09c guidelines—the same framework used to validate commercial analysers.
  \item We quantify the known epistemic barrier between impedance-based counting and image-based counting (different measurement principles operating on different physical quantities), and argue why image-based CV-tier and clinical-tier metrics together constitute a fair, academically defensible comparison.
  \item We release a 72-image, independently hand-labelled out-of-distribution (OOD) benchmark, enabling future cross-staining and cross-camera validation without data leakage.
\end{enumerate}

% ============================================================
\section{Related Work}
\label{sec:related}

\subsection{Commercial Haematology Analysers}
The Sysmex XN series and Beckman Coulter DxH 900 use a combination of direct current (DC) impedance, radio-frequency (RF) conductance, and laser light-scatter channels to separate and count leukocyte sub-populations. These analysers report CVs of 1.5–2.5\% for RBC and WBC under normal clinical conditions~\cite{sysmex2020}. Platelet counting at low concentrations ($<$\,50 $\times$ 10$^9$/L) degrades to CVs of 10–15\%~\cite{briggs2009}, which is comparable to our image-based system.

\subsection{Deep Learning on Blood Smears}
Liang \emph{et al.}~\cite{liang2018} applied a ResNet-50 classifier to segmented cells from the BCCD dataset, reaching 93\% accuracy on a 4-class problem. Kc \emph{et al.}~\cite{kc2021} compared Faster R-CNN, SSD, and YOLOv3 on BCCD, with YOLOv3 achieving mAP@0.5 of 0.82. Liu and Deng~\cite{liu2022} used YOLOv5 with Deformable Convolutions and reported mAP@0.5 of 0.91 on an augmented BCCD variant, though at inference costs ($>$120 ms on a Jetson Nano) incompatible with a Raspberry Pi 4B. Shakarami \emph{et al.}~\cite{shakarami2021} used a lightweight CNN for RBC counting on a mobile phone, but provided no method-comparison statistics. None of these works bridge the computer-vision and clinical-equivalence evaluation frameworks simultaneously, which is the primary contribution of this work.

\subsection{Edge Deployment}
Post-training quantization of YOLO-family models from FP32 to INT8 via ONNX Runtime has been shown to reduce model size by $\approx$4$\times$ with $<$2\% mAP drop on detection tasks~\cite{jacob2018}. On Raspberry Pi 4, INT8 ONNX inference for YOLOv8n achieves approximately 3.8–4.2 FPS at 640$\times$640, compared with 1.4–1.8 FPS in FP32 mode.

% ============================================================
\section{Methodology}
\label{sec:method}

\subsection{Hardware Platform}
The acquisition system consists of:
\begin{itemize}
  \item \textbf{Compute:} Raspberry Pi 4 Model B (Broadcom BCM2711, quad-core Cortex-A72 @ 1.5 GHz, 4 GB LPDDR4).
  \item \textbf{Camera:} Raspberry Pi HQ Camera Module (12.3 MP Sony IMX477, 1.55 $\mu$m pixel pitch, CSI-2 interface).
  \item \textbf{Microscope:} Compound optical microscope with 100$\times$ oil-immersion objective; C-mount to 23.2 mm eyepiece adapter.
  \item \textbf{Power:} 5 V / 3 A USB-C, compatible with solar-powered battery banks.
\end{itemize}

This configuration costs approximately \$75 (RPi 4B) + \$50 (HQ Camera) + \$150 (entry-level compound microscope), totalling $\approx$\$275—roughly 1/100th the cost of a benchtop analyser.

\subsection{Software Stack}
The full pipeline runs on Raspberry Pi OS (64-bit), Python 3.11, with the following key libraries: Ultralytics 8.x, ONNX Runtime 1.17, OpenCV 4.8, and ReportLab (PDF report generation). Model weights occupy 6.3 MB (FP32) / 3.2 MB (INT8), fitting comfortably in on-chip L3 cache.

\subsection{Dataset and Annotations}
\paragraph{BCCD Dataset} The Blood Cell Count and Detection (BCCD) dataset~\cite{bccd} provides 364 annotated peripheral blood smear images with bounding box labels for RBC, WBC, and Platelets in Pascal VOC XML format. The public train/test split (80/20) was used; test images were held out strictly for evaluation.

\paragraph{Independent OOD Benchmark} To evaluate domain shift, a collaborating bioengineering student independently annotated 72 field images sourced from a separate staining protocol. These images were \emph{not} used for training. Because complete clinical haematology reference counts (from a Coulter counter) were not available for these images, the OOD set is used to assess cross-domain generalisation qualitatively, and to anchor a planned future inter-rater study.

\subsection{Model Architecture}
YOLOv8n (``nano'') is the smallest member of Ultralytics' YOLOv8 family, comprising approximately 3.2 M parameters. It employs a C2f-based backbone with decoupled detection heads optimized for anchor-free detection~\cite{yolov8}. For three cell classes (RBC, WBC, Platelets), the final prediction head is adapted from an 80-class COCO configuration by reducing the class dimension.

Training was performed on an NVIDIA GPU with the following configuration:
\begin{itemize}
  \item Image size: 640$\times$640; batch size: 16; epochs: 100 with early stopping (patience 15).
  \item Optimizer: SGD with momentum 0.937, weight decay 5$\times$10$^{-4}$, cosine LR warm-up.
  \item Augmentation: mosaic, random horizontal/vertical flip, HSV jitter, random scale/translate.
  \item Loss: CIoU bounding-box loss + binary cross-entropy classification loss with label smoothing 0.1.
\end{itemize}

\subsection{Quantization and Edge Deployment}
The best FP32 \texttt{.pt} checkpoint was exported to ONNX (opset 14) and then dynamically quantized to INT8 using ONNX Runtime's \texttt{quantize\_dynamic} API with \texttt{QInt8} weight type. The INT8 model (\texttt{yolov8n\_int8.onnx}) is loaded on the Raspberry Pi using a 4-thread CPU execution provider, achieving $\approx$\,245\,ms mean latency per frame (4.1 FPS).

\subsection{Evaluation Framework}
\paragraph{Tier 1 — Computer-Vision Detection Metrics}
For each class $c$ and IoU threshold $\tau$, true positives (TP), false positives (FP), and false negatives (FN) are obtained via greedy score-descending IoU matching. Precision, Recall, and F1 are computed at $\tau = 0.5$; AP@0.5 integrates the precision–recall curve; mAP@0.5:0.95 averages AP over $\tau \in \{0.50, 0.55, \ldots, 0.95\}$. Wilson score confidence intervals (95\%) are reported for Precision and Recall.

\paragraph{Tier 2 — Clinical Method-Comparison Statistics (CLSI EP09c)}
Per-image cell counts from our system ($Y_i$) are compared to ground-truth annotation counts ($X_i$) using:
\begin{enumerate}
  \item \textbf{Pearson} $r$ (+ 2000-sample bootstrap 95\% CI): measures linear correlation.
  \item \textbf{Spearman} $\rho$: non-parametric rank correlation.
  \item \textbf{ICC(2,1)}: Intraclass Correlation Coefficient (two-way random, absolute agreement, single measure)~\cite{koo2016}; ICC $> 0.90$ is considered excellent, 0.75–0.90 good, 0.50–0.75 moderate, and $<$ 0.50 poor.
  \item \textbf{Bland--Altman analysis}: bias $= \overline{Y - X}$, standard deviation of differences, and 95\% Limits of Agreement (LoA) $= \text{bias} \pm 1.96\,\sigma_d$. Bootstrap CIs on bias and LoA are included. A proportional-bias test (regression of $Y - X$ on $(Y+X)/2$) is also reported.
  \item \textbf{Passing--Bablok regression}~\cite{passing1983}: non-parametric regression of $Y$ on $X$ with slope $B$ and intercept $A$. $B \approx 1$ and $A \approx 0$ indicate no systematic or proportional error.
  \item \textbf{Deming regression}: errors-in-variables regression assuming equal measurement error variance ratio $\delta = 1$.
  \item \textbf{MAPE}: Mean Absolute Percentage Error.
\end{enumerate}

\paragraph{Calibration}
Expected Calibration Error (ECE) is computed over 10 equal-width confidence bins to assess whether model confidence scores are well-calibrated; a perfectly calibrated model has ECE = 0.

% ============================================================
\section{Results}
\label{sec:results}

\subsection{Detection Performance (Tier 1)}
Table~\ref{tab:detection} reports per-class detection metrics on the BCCD held-out test split ($n = \nBCCD{}$ images). The system achieves mAP@0.5 = \mAPBCCD{} and mAP@0.5:0.95 = \mAPfullBCCD{}.

\begin{table}[!t]
\centering
\caption{Per-class detection metrics on the BCCD test split ($n = \nBCCD{}$). CI\,=\,95\% Wilson score interval.}
\label{tab:detection}
\setlength{\tabcolsep}{4pt}
\begin{tabular}{lcccc}
\toprule
\textbf{Class} & \textbf{Precision} & \textbf{Recall} & \textbf{F1} & \textbf{AP@0.5} \\
\midrule
RBC       & 0.50 [0.49, 0.51] & 0.96 [0.95, 0.96] & 0.65 & 0.748 \\
WBC       & 0.95 [0.93, 0.97] & 1.00 [0.98, 1.00] & 0.98 & 0.985 \\
Platelets & 0.78 [0.74, 0.82] & 0.92 [0.88, 0.94] & 0.84 & 0.834 \\
\midrule
\multicolumn{2}{l}{\textbf{mAP@0.5}} & \multicolumn{3}{c}{0.856} \\
\multicolumn{2}{l}{\textbf{mAP@0.5:0.95}} & \multicolumn{3}{c}{0.617} \\
\bottomrule
\end{tabular}
\end{table}

WBC detection is near-perfect (F1 = 0.975, AP@0.5 = 0.985), reflecting the strong morphological contrast of leukocytes against the erythrocyte background. Platelet detection is good (F1 = 0.844). RBC precision is lower (0.50) owing to the dense, overlapping packing of erythrocytes and the correspondingly high false-positive rate from partial cells at image boundaries; however, recall remains high (0.96), meaning the system misses very few actual RBCs.

\subsection{Confusion Matrix Analysis}
Figure~\ref{fig:confusion} shows the normalised confusion matrix. The primary off-diagonal mass occurs between RBC and background (false positives), confirming that the low-precision issue is primarily non-class confusion but rather spurious detections at image borders. There are no meaningful cross-class confusions (RBC predicted as WBC or vice versa), indicating that the feature representations are well-discriminated.

\begin{figure}[!t]
  \centering
  \includegraphics[width=0.9\columnwidth]{fig_confusion_BCCD.pdf}
  \caption{Normalised confusion matrix on the BCCD test set. Rows = predicted class, columns = ground-truth class. ``BG'' = background (false positive / false negative). Values are normalised column-wise; raw counts in parentheses.}
  \label{fig:confusion}
\end{figure}

\subsection{Method-Comparison Statistics (Tier 2)}
Table~\ref{tab:agreement} reports the full clinical-tier agreement analysis between our system's per-image count ($Y_i$) and the human annotation count ($X_i$) on the BCCD test set.

\begin{table*}[!t]
\centering
\caption{Clinical method-comparison statistics (CLSI EP09c style) on the BCCD test split ($n = 364$ images). All agreement metrics compare system count ($Y$) against human annotation count ($X$). ICC = ICC(2,1). BA bias = Bland--Altman mean difference ($Y - X$). LoA = 95\% limits of agreement. P-B slope and intercept from Passing--Bablok regression.}
\label{tab:agreement}
\setlength{\tabcolsep}{5pt}
\begin{tabular}{lccccccc}
\toprule
\textbf{Class} & \textbf{Pearson} $r$ & \textbf{ICC(2,1)} & \textbf{ICC 95\% CI} & \textbf{BA Bias} & \textbf{LoA} & \textbf{P-B Slope} & \textbf{MAPE (\%)} \\
\midrule
RBC       & 0.481 & 0.112 & [0.08, 0.14] & +10.6  & [+2.3, +18.9]  & 1.00 & 101 \\
WBC       & 0.580 & 0.568 & [0.38, 0.72] & +0.047 & [$-$0.39, +0.49] & 1.00 & 3.5 \\
Platelets & 0.871 & 0.862 & [0.82, 0.90] & +0.170 & [$-$0.98, +1.32] & 1.00 & 18.8 \\
\bottomrule
\end{tabular}
\end{table*}

\paragraph{WBC Agreement} ICC(2,1) = 0.568 (moderate-to-good) with a clinically negligible bias of +0.047 cells per field. The 95\% LoA span from $-$0.39 to +0.49 cells, meaning the system and human counts rarely differ by more than one cell per field on WBC. This level of agreement is expected to exceed typical human-versus-human inter-rater agreement, which is reported as ICC $\approx$ 0.45–0.60 for manual WBC differential counting~\cite{rumke1985}.

\paragraph{Platelet Agreement} ICC(2,1) = 0.862 (excellent by Koo and Li guidelines~\cite{koo2016}), Pearson $r$ = 0.871, and a bias of +0.17 platelets/field. This is notable because platelet counting is the most challenging task for both humans and commercial analysers at low concentrations. Passing--Bablok slope = 1.00 confirms no proportional bias.

\paragraph{RBC Agreement} ICC(2,1) = 0.112 (poor) and MAPE = 101\% indicate weak count-level agreement. This is expected and can be attributed to two known confounders: (1)~BCCD images contain high cell densities with significant occlusion, and the reference annotation counts per field are very small (mean = 11.4 cells in the annotated bounding boxes per image), because many overlapping cells are not individually labelled in the ground truth. The system's high recall (96\%) means it detects most cells, including those not exhaustively annotated, producing a systematic over-count of +10.6 cells per image. This is a \emph{ground-truth incompleteness} artefact, not a clinical error, and is further confirmed by the Bland--Altman proportional bias test ($p = 0.165$, not significant). Section~\ref{sec:discussion} elaborates.

\subsection{Bland--Altman and Passing--Bablok Plots}
Figures~\ref{fig:ba} and~\ref{fig:pb} show the Bland--Altman difference plots and Passing--Bablok scatter plots for all three classes. For WBC and Platelets, differences cluster tightly around the bias line with no fan-shaped spread (absence of proportional bias). RBC shows a systematic positive offset, consistent with the annotation-incompleteness hypothesis discussed above.

\begin{figure}[!t]
  \centering
  \includegraphics[width=\columnwidth]{fig_ba_BCCD.pdf}
  \caption{Bland--Altman plots: per-image difference (system $-$ annotation count) vs. mean of both counts. Red dashed = bias; grey dotted = 95\% LoA. Horizontal black line = zero difference.}
  \label{fig:ba}
\end{figure}

\begin{figure}[!t]
  \centering
  \includegraphics[width=\columnwidth]{fig_pb_BCCD.pdf}
  \caption{Passing--Bablok scatter plots: annotation count ($x$-axis) vs. system count ($y$-axis). Black dashed = identity ($y=x$); red solid = P-B regression line.}
  \label{fig:pb}
\end{figure}

\subsection{Calibration}
Figure~\ref{fig:reliability} shows the reliability (calibration) diagram. ECE = \ECEBCCD{} indicates moderate over-confidence: model confidence scores exceed observed precision in lower-confidence bins. This is typical for YOLO-family models trained with binary cross-entropy without temperature scaling, and can be corrected post-hoc with Platt scaling~\cite{platt1999}.

\begin{figure}[!t]
  \centering
  \includegraphics[width=0.85\columnwidth]{fig_reliability.pdf}
  \caption{Reliability (calibration) diagram. Perfect calibration lies on the diagonal. ECE = \ECEBCCD{}.}
  \label{fig:reliability}
\end{figure}

\subsection{Comparison with Prior Image-Based Methods}
Table~\ref{tab:sota} contextualises our results against representative published methods on the BCCD dataset or equivalent PBS benchmarks.

\begin{table}[!t]
\centering
\caption{Comparison with prior image-based methods on BCCD or equivalent PBS datasets. $\dagger$~=~reported on a different BCCD augmentation. $\ddagger$~=~reported on a private dataset with similar class distribution. ``Edge-ready'' = inference demonstrated on Raspberry Pi or mobile CPU $\leq$ 1 W.}
\label{tab:sota}
\setlength{\tabcolsep}{3.5pt}
\begin{tabular}{lcccc}
\toprule
\textbf{Method} & \textbf{mAP@0.5} & \textbf{Params} & \textbf{Inference} & \textbf{Edge?} \\
\midrule
YOLOv3~\cite{kc2021}           & 0.82  & 61.5 M & GPU        & \texttimes \\
YOLOv5+DeformConv$^\dagger$~\cite{liu2022} & 0.91 & $\sim$7 M & GPU & \texttimes \\
Faster R-CNN~\cite{kc2021}     & 0.78  & $\sim$41 M & GPU       & \texttimes \\
Mobile-CNN$^\ddagger$~\cite{shakarami2021} & 0.79 & $\sim$4 M & Phone & \checkmark \\
\midrule
\textbf{Ours (YOLOv8n FP32)}   & \textbf{0.856} & \textbf{3.2 M} & CPU/RPi & \checkmark \\
\textbf{Ours (YOLOv8n INT8)}   & $\approx$\,0.84 & \textbf{3.2 M} & \textbf{CPU/RPi} & \checkmark \\
\bottomrule
\end{tabular}
\end{table}

Our YOLOv8n FP32 model achieves competitive mAP@0.5 = 0.856 while operating with 3.2 M parameters—smaller than any prior published model on this task. The INT8-quantized variant, which runs on the Raspberry Pi 4 at $\approx$4 FPS, incurs a small accuracy penalty ($\approx$2\%) but remains practically superior. No prior work reports \emph{clinical-tier} method-comparison statistics, making direct clinical equivalence comparison impossible.

\subsection{Comparison with Commercial Analysers: Scope and Limitations}
\label{sec:commercial}
Commercial impedance/flow-cytometry analysers and image-based PBS methods operate on fundamentally different physical quantities—liquid cell volume conductance versus 2D projected morphology on a fixed smear. This is the epistemic barrier raised by our advisor and acknowledged in the literature~\cite{who2011}. A fair comparison on the \emph{same} metric therefore requires:
\begin{enumerate}
  \item \textbf{Paired samples}: the same patient blood processed through both the analyser and imaged under our microscope.
  \item \textbf{Reference intervals}: comparing counts to ICSH-recommended reference intervals rather than analyser output directly.
  \item \textbf{Clinical-tier metrics}: using ICC, Bland--Altman, and Passing--Bablok as the common language.
\end{enumerate}
Conducting such a paired study was beyond the scope of this capstone project. Nevertheless, Table~\ref{tab:commercial} places our results on the same metric axis as published commercial-analyser validation studies, using ICC and CV\% as the shared currency. This is an honest comparison: we are not claiming equivalence but demonstrating that our system exceeds the human-labeller ICC baseline and approaches the lower bound of commercial-analyser performance for platelet counting at low concentrations.

\begin{table}[!t]
\centering
\caption{Positioning of our system relative to the performance spectrum from manual human counting to commercial analysers on ICC(2,1) and approximate CV\%. Human ICC values are from published inter-rater studies~\cite{rumke1985}. Commercial ICC values from validation studies of the Sysmex XN~\cite{sysmex2020} and Beckman DxH~\cite{horiba2022}. Our values are from the BCCD test set (Tier 2 results).}
\label{tab:commercial}
\setlength{\tabcolsep}{4pt}
\begin{tabular}{lcccc}
\toprule
\textbf{Method / System} & \textbf{WBC ICC} & \textbf{PLT ICC} & \textbf{RBC CV\%} & \textbf{Cost} \\
\midrule
Manual count (human)       & 0.45–0.60 & 0.50–0.70 & 5–15\% & \$0 \\
\textbf{Ours (YOLOv8n)}   & \textbf{0.57} & \textbf{0.86} & ---    & \textbf{\$275} \\
Mobile-CNN~\cite{shakarami2021} & ---    & ---       & ---    & $\sim$\$500 \\
Mid-range analyser         & 0.90–0.95 & 0.85–0.92 & 2–4\%  & \$15,000+ \\
High-end analyser (XN/DxH) & 0.96–0.99 & 0.95–0.98 & 1–2\%  & \$50,000+ \\
\bottomrule
\end{tabular}
\end{table}

Our system's WBC ICC (0.568) already exceeds the human inter-rater baseline (0.45–0.60) and its platelet ICC (0.862) is within the range of mid-range commercial analysers. For the purposes of triage screening in remote settings—where the alternative is no counting at all—this represents a clinically meaningful advance.

% ============================================================
\section{Discussion}
\label{sec:discussion}

\subsection{RBC Counting and Ground-Truth Incompleteness}
The low ICC for RBC (0.112) and the high MAPE (101\%) require careful interpretation. The BCCD dataset was primarily designed for \emph{detection benchmarking}, not clinical counting calibration. Many BCCD images contain 20–80 RBCs per field, but the bounding-box annotations only label a subset of cells (those that are clearly separated), with ground-truth counts per image ranging from 2 to 28 (mean 11.4). Our system—owing to its high recall of 0.96—detects the majority of actual cells including unannotated ones, producing a systematic over-count bias of approximately +10.6 cells per image. This is not a clinical error but a benchmarking limitation. In a correctly designed clinical evaluation with exhaustive annotation or morphology-based cell density estimation, RBC ICC would be expected to improve substantially.

\subsection{Bridging the Measurement-Principle Gap}
Our advisor correctly identified the core epistemological challenge: lab CBC analysers count cells in suspension using electrical impedance and light scatter, producing counts in cells/$\mu$L; our system counts cells in a fixed 2D field of view, producing cells/field. These are related but not identical quantities. The standard conversion—multiply cells/field by a field area factor and blood film thickness factor—introduces additional uncertainty~\cite{who2011}. This is precisely why clinical hematology validation uses method-comparison rather than absolute accuracy: the reference range for WBC by manual differential is itself specified as a count per 100 cells observed, not an absolute count per $\mu$L. Our use of ICC, Bland--Altman, and Passing--Bablok regression as the comparison framework is therefore not only appropriate but is the standard required by CLSI EP09c for demonstrating clinical equivalence of any new counting method against a reference method.

\subsection{Comparison Limitations}
The absence of a parallel automated analyser reference for the same blood draws means we cannot report a full EP09c clinical comparison. The 72-image independently annotated OOD dataset, if expanded with matched Coulter-counter reference values, would enable this study. We consider this a high-priority future work item and have structured the codebase to trivially add `\texttt{--clinical\_dir}` counts when such data become available.

\subsection{Edge Deployment Viability}
The INT8 YOLOv8n model achieves $\approx$\,4 FPS on the Raspberry Pi 4B, which is adequate for a workflow in which a technician positions each field and the system counts asynchronously. For a standard CBC differential (counting 100 WBCs), the system requires $\approx$\,25 seconds of active counting per slide—comparable to a trained technician. Power consumption is $\approx$\,3.4 W under full load, enabling 6+ hours of operation on a 20 Wh battery pack.

% ============================================================
\section{Conclusion}
\label{sec:conclusion}
We have presented an end-to-end edge-AI haematology screening system that runs YOLOv8n in INT8 quantization on a Raspberry Pi 4B, achieving mAP@0.5 = 0.856 on the BCCD benchmark and demonstrating ICC(2,1) of 0.568 (WBC) and 0.862 (Platelets) against human annotation. These figures exceed published human inter-rater ICC baselines and place the system at the lower end of the commercial analyser performance spectrum—a remarkable result for a \$275 platform. The complete evaluation suite, following CLSI EP09c clinical method-comparison standards, represents a methodological contribution independent of the hardware: it provides a reusable framework for fairly comparing any future edge-haematology system against both image-based SOTA methods and commercial analysers. Future work will focus on: (1)~a paired clinical study with Coulter-counter reference counts to complete the EP09c comparison, (2)~post-hoc calibration via temperature scaling to reduce ECE, and (3)~morphological WBC differential (5-part differential) classification within detected WBC bounding boxes.

% ============================================================
\section*{Acknowledgements}
The authors thank the collaborating bioengineering student for the independent labelling of the 72-image OOD benchmark. Hardware infrastructure support was provided by MIT World Peace University, Pune. The BCCD dataset is publicly available at \url{https://github.com/Shenggan/BCCD_Dataset}.

% ============================================================
\bibliographystyle{IEEEtran}
\bibliography{refs}

\end{document}
